\documentclass[10pt,twocolumn]{article}
\usepackage[a4paper,top=1.9cm,bottom=2.2cm,left=1.7cm,right=1.7cm,columnsep=0.6cm]{geometry}
\usepackage{times}
\usepackage[T1]{fontenc}
\usepackage{graphicx}
\usepackage{booktabs}
\usepackage{caption}
\usepackage{amsmath}
\usepackage{hyperref}
\usepackage{titlesec}
\usepackage{enumitem}
\usepackage{abstract}
\usepackage{fancyhdr}

\graphicspath{{../figures/}}

\hypersetup{colorlinks=true, linkcolor=black, citecolor=black, urlcolor=blue}

\titleformat{\section}{\normalfont\large\bfseries}{\Roman{section}.}{0.6em}{}
\titleformat{\subsection}{\normalfont\bfseries}{\Alph{subsection}.}{0.5em}{}
\renewcommand{\abstractnamefont}{\normalfont\bfseries}
\renewcommand{\abstracttextfont}{\normalfont\itshape}
\setlength{\absleftindent}{0.5cm}
\setlength{\absrightindent}{0.5cm}

\pagestyle{fancy}
\fancyhf{}
\renewcommand{\headrulewidth}{0pt}
\fancyfoot[C]{\thepage}

\captionsetup[figure]{font=small,labelfont=bf}
\captionsetup[table]{font=small,labelfont=bf}

\title{\LARGE\bfseries Influence, Brokerage, and Community Structure\\ in the 2004 U.S. Political Blogosphere}

\author{
  \large Student Name\\[2pt]
  \normalsize Student ID -- SE4095 Social Network Theory and Practice, Assignment 01
}

\date{}

\begin{document}
\twocolumn[
  \begin{@twocolumnfalse}
    \maketitle
    \begin{abstract}
    This report analyses the political blogosphere hyperlink network compiled by Adamic and Glance
    (2005) to investigate how influence and brokerage relate to political homophily. The raw graph
    (1490 nodes, 19{,}090 directed hyperlinks) is cleaned into a simple directed graph and restricted,
    for all distance-based and structural calculations, to its largest weakly connected component
    (1222 nodes, 19{,}021 edges). PageRank and betweenness centrality are computed to distinguish
    cited authority from structural brokerage; the two rankings correlate strongly (Spearman
    $\rho=0.76$) but agree on only a third of their top-10 nodes. A controlled removal experiment
    shows that deleting the top 3\% of blogs by betweenness fragments the network into roughly six
    times as many Louvain communities and sharply reduces agreement with the known political labels,
    while an equal-size random removal leaves the community structure almost unchanged. Louvain
    modularity maximisation, run without ever seeing the ground-truth labels, recovers the liberal
    and conservative divide as two large, $>90\%$-pure communities, evaluated here with a pairwise
    confusion matrix (accuracy 0.88, F1 0.87 on the original graph) that avoids the community-ID
    permutation problem entirely.
    \end{abstract}
    \vspace{0.6cm}
  \end{@twocolumnfalse}
]

\section{Introduction}
The 2005 Political Blogs dataset~\cite{adamic2005} records hyperlinks among 1490 U.S. political
blogs crawled around the 2004 presidential election, together with each blog's political
orientation (liberal or conservative). This report investigates three questions: (1) which blogs
act as authorities versus brokers in the hyperlink structure; (2) whether removing the network's
structural brokers destabilises its community structure more than removing an equal number of
random blogs; and (3) whether an unsupervised, label-free community-detection procedure recovers
the known political split. Ground-truth orientation labels are preserved throughout but are used
only in the final evaluation stage, never to construct, tune, or select communities.

\section{Dataset and Graph Construction}
\label{sec:data}

\subsection{Representation}
The GML file is loaded as a directed graph: a hyperlink from blog $u$ to blog $v$ is not evidence
that $v$ links back, and direction is required by both centrality measures used later. Edges are
treated as unweighted -- the file provides no independently verifiable count of repeated links, so
edge multiplicity is collapsed into a single boolean ``linked'' relationship rather than used as an
invented weight.

\subsection{Preprocessing log}
As distributed, \texttt{polblogs.gml} is a directed \emph{multigraph}, which the standard NetworkX
GML reader rejects unless \texttt{multigraph~1} is declared in the header -- adding that declaration
was the first preprocessing action. Table~\ref{tab:preprocess} documents every subsequent decision
and the resulting graph size at each stage.

\begin{table}[h]
\centering
\caption{Preprocessing log (raw file $\rightarrow$ analysis graph)}
\label{tab:preprocess}
\small
\begin{tabular}{@{}lr@{}}
\toprule
Step & Count \\
\midrule
Raw nodes / edges (multigraph) & 1490 / 19{,}090 \\
Self-loops removed & 3 \\
Duplicate (source,target) pairs collapsed & 65 \\
Nodes missing orientation label & 0 \\
Cleaned simple graph (nodes / edges) & 1490 / 19{,}022 \\
Isolated nodes & 266 \\
Weakly connected components & 268 \\
Largest WCC size & 1222 \\
Strongly connected components & 688 \\
Largest SCC size & 793 \\
\textbf{Analysis graph (largest WCC)} & \textbf{1222 / 19{,}021} \\
\bottomrule
\end{tabular}
\end{table}

\subsection{Component choice for distance-based measures}
The largest strongly connected component covers only 793 of 1490 nodes (53\%), which would discard
almost half the network. The largest \emph{weakly} connected component (1222 nodes, 65\% of the
cleaned graph) is used instead for every centrality, community-detection, and pairwise-evaluation
calculation in this report: it retains every blog reachable through the network in some direction,
matching the notion of ``part of the same information-flow structure'' better than requiring mutual
reachability. The 268 excluded nodes are 266 true isolates plus one disconnected pair.

\subsection{Graph properties}
Table~\ref{tab:props} reports density and degree statistics for the full cleaned graph and the
1222-node analysis graph used from Section~\ref{sec:centrality} onward.

\begin{table}[h]
\centering
\caption{Graph properties}
\label{tab:props}
\small
\begin{tabular}{@{}lrr@{}}
\toprule
Property & Full cleaned & Analysis (LWCC) \\
\midrule
Nodes & 1490 & 1222 \\
Edges & 19{,}022 & 19{,}021 \\
Density & 0.0086 & 0.0127 \\
Weakly conn. components & 268 & 1 \\
Strongly conn. components & 688 & 420 \\
Largest SCC fraction & 0.53 & 0.65 \\
Mean in/out-degree & 12.77 & 15.57 \\
Median in-degree & 2 & 3 \\
Median out-degree & 4 & 7 \\
Max in-degree & 337 & 337 \\
Max out-degree & 256 & 256 \\
\bottomrule
\end{tabular}
\end{table}

\begin{figure}[h]
\centering
\includegraphics[width=0.98\linewidth]{degree_distribution_analysis_LWCC.png}
\caption{In- and out-degree distributions of the analysis graph (log-log). Both are heavy-tailed: a
small number of hub blogs receive the majority of in-links.}
\label{fig:degdist}
\end{figure}

\begin{figure}[h]
\centering
\includegraphics[width=0.98\linewidth]{network_overview_by_orientation.png}
\caption{Analysis graph coloured by the known (but unused-in-detection) political orientation.
Colour is descriptive only; no calculation in this report uses it before Section~\ref{sec:eval}.}
\label{fig:overview}
\end{figure}

\subsection{Interpretation}
Figure~\ref{fig:overview} shows two dense, largely separate clusters connected by comparatively few
cross-links -- a clear visual signature of homophily: blogs overwhelmingly link within their own
political side. Figure~\ref{fig:degdist} shows heavy-tailed in- and out-degree distributions
consistent with a small number of highly visible blogs dominating attention while most blogs are
linked to only once or twice.

\section{Centrality Selection and Analysis}
\label{sec:centrality}

\subsection{Chosen measures and justification}
Two taught measures are selected: \textbf{PageRank} and \textbf{betweenness centrality}. PageRank is
the only taught measure built specifically for directed link graphs with divided outgoing influence:
a link from a blog that cites 300 others counts far less than one from a blog that cites five, a
distinction plain in-degree cannot make. Betweenness is the only taught measure that scores
brokerage on shortest paths rather than direct or link-weighted importance -- exactly the property
needed to design the removal experiment in Section~\ref{sec:experiment}. Both are computed on the
directed analysis graph with the standard damping factor $\alpha=0.85$ for PageRank (tight
convergence, $\mathrm{tol}=10^{-10}$) and directed shortest paths for betweenness, where unreachable
pairs contribute zero, the standard convention for a graph that is only weakly connected.

\begin{table}[h]
\centering
\caption{Top-10 blogs by PageRank}
\label{tab:pr}
\small
\begin{tabular}{@{}rlll@{}}
\toprule
Rank & Blog & Orientation & PageRank \\
\midrule
1 & dailykos.com & liberal & 0.0189 \\
2 & atrios.blogspot.com & liberal & 0.0160 \\
3 & instapundit.com & conservative & 0.0133 \\
4 & blogsforbush.com & conservative & 0.0132 \\
5 & talkingpointsmemo.com & liberal & 0.0131 \\
6 & michellemalkin.com & conservative & 0.0115 \\
7 & drudgereport.com & conservative & 0.0113 \\
8 & washingtonmonthly.com & liberal & 0.0111 \\
9 & powerlineblog.com & conservative & 0.0094 \\
10 & andrewsullivan.com & conservative & 0.0091 \\
\bottomrule
\end{tabular}
\end{table}

\begin{table}[h]
\centering
\caption{Top-10 blogs by betweenness centrality}
\label{tab:bc}
\small
\begin{tabular}{@{}rlll@{}}
\toprule
Rank & Blog & Orientation & Betweenness \\
\midrule
1 & blogsforbush.com & conservative & 0.1467 \\
2 & atrios.blogspot.com & liberal & 0.0611 \\
3 & instapundit.com & conservative & 0.0512 \\
4 & dailykos.com & liberal & 0.0369 \\
5 & newleftblogs.blogspot.com & liberal & 0.0308 \\
6 & madkane.com/notable.html & liberal & 0.0302 \\
7 & wizbangblog.com & conservative & 0.0273 \\
8 & lashawnbarber.com & conservative & 0.0243 \\
9 & hughhewitt.com & conservative & 0.0230 \\
10 & washingtonmonthly.com & liberal & 0.0219 \\
\bottomrule
\end{tabular}
\end{table}

\subsection{Ranking comparison}
Spearman's rank correlation between the full PageRank and betweenness rankings is
$\rho = 0.76$ ($p \approx 3\times10^{-228}$) -- a strong positive association, since a blog that
receives many selective in-links is also likely to lie on many shortest paths. The two rankings are
nonetheless distinct: only 5 of the top 10 nodes are shared (Jaccard $=0.33$).

\subsection{Node interpretation}
\textbf{dailykos.com} ranks \#1 by PageRank but only \#4 by betweenness: it is the most-cited
liberal hub, but so many blogs link to it directly that fewer shortest paths need to route
\emph{through} it. \textbf{blogsforbush.com} is \#4 by PageRank but \#1 by betweenness by a wide
margin (0.147, more than double the runner-up) -- it is not the most-cited conservative blog, but it
appears to be the main crossing point that shortest paths between many liberal and conservative
blogs are forced through. \textbf{atrios.blogspot.com} ranks highly on both measures (\#2 and \#2),
combining cited authority with structural bridging, a rarer and arguably more structurally important
profile than either extreme alone. A high score on either measure describes a blog's position in the
hyperlink graph, not necessarily its real-world political influence.

\section{Centrality-Informed Community Experiment}
\label{sec:experiment}

\subsection{Design}
Betweenness centrality is used as the intervention measure, since Section~\ref{sec:centrality}
identifies a small set of blogs -- led by blogsforbush.com -- that structurally bridge the two
political sides. If these blogs are genuinely load-bearing brokers, removing them should fragment
community structure more than removing an equal number of random blogs. The top 3\% of nodes by
betweenness (37 of 1222) are removed to build a \emph{targeted} graph, large enough to meaningfully
stress the network while leaving 97\% intact; a \emph{random} graph removes the same count of nodes
chosen uniformly at random with a fixed seed (Python \texttt{random.Random(42)}). Both removal
graphs (1185 nodes each) are evaluated only on their own remaining nodes.

\section{Community Detection}
\label{sec:community}

\subsection{Method}
Louvain modularity maximisation~\cite{blondel2008} is used. Louvain assumes disjoint, reasonably
dense communities and optimises a modularity objective, with no assumption of balanced sizes,
hierarchy, or overlap; it is defined for undirected graphs. The directed analysis graph is converted
with an edge kept if a hyperlink exists in either direction -- a deliberate choice, since the
community question is about mutual reference rather than the direction of any single link, and
direction was already fully exploited at the centrality stage. Louvain's node-visiting order is
randomised, so five seeds (0--4) are run for each of the three graphs from Section~\ref{sec:experiment};
seed 0 is the declared run used for the visualisations and profile table, and mean/standard
deviation across seeds are reported for every metric.

\subsection{Results across the three graphs}
Table~\ref{tab:summary} summarises modularity and the external metrics (defined in
Section~\ref{sec:eval}) across five Louvain seeds for the original graph and both removal
conditions.

\begin{table}[h]
\centering
\caption{Mean $\pm$ std over 5 Louvain seeds}
\label{tab:summary}
\small
\begin{tabular}{@{}lrrr@{}}
\toprule
 & Original & Targeted rem. & Random rem. \\
\midrule
$N$ & 1222 & 1185 & 1185 \\
\# communities & $10.2\pm1.3$ & $61.0\pm1.2$ & $13.0\pm1.2$ \\
Modularity & $0.427\pm0.000$ & $0.433\pm0.001$ & $0.423\pm0.001$ \\
Accuracy & $0.883\pm0.003$ & $0.764\pm0.029$ & $0.878\pm0.002$ \\
Precision & $0.919\pm0.007$ & $0.912\pm0.009$ & $0.913\pm0.005$ \\
Recall & $0.842\pm0.003$ & $0.585\pm0.057$ & $0.835\pm0.002$ \\
F1 & $0.878\pm0.002$ & $0.711\pm0.043$ & $0.872\pm0.002$ \\
\bottomrule
\end{tabular}
\end{table}

Removing the 37 highest-betweenness blogs raises the community count roughly six-fold (10--11 to
around 60) and collapses recall from $\approx0.84$ to $0.55$--$0.68$: many same-orientation blogs end
up scattered across small fragments rather than one large community. Removing 37 random blogs barely
moves any metric. This gap is direct evidence that the highest-betweenness blogs are structurally
load-bearing: losing them breaks the paths holding each political side together as one community,
whereas losing an equal number of typical blogs does not. Modularity rises slightly under targeted
removal even as label agreement collapses -- removing cross-cutting brokers leaves many small,
internally dense fragments, which modularity rewards even though they are less politically coherent
as groups than the two original giant communities.

\begin{figure}[h]
\centering
\includegraphics[width=0.98\linewidth]{communities_original.png}
\caption{Declared-run (seed 0) Louvain communities, original graph. Two giant communities dominate.}
\label{fig:comm_orig}
\end{figure}

\begin{figure}[h]
\centering
\includegraphics[width=0.98\linewidth]{communities_targeted_removal.png}
\caption{Communities after removing the top 37 blogs by betweenness: the core fragments into many
small pieces and a ring of newly-disconnected peripheral nodes appears.}
\label{fig:comm_targeted}
\end{figure}

\subsection{Community profile (declared run, original graph)}
Table~\ref{tab:profile} profiles the largest communities using the two Part~B centralities. The two
dominant communities (649 and 518 blogs) are highly pure (93.8\% conservative, 97.7\% liberal) and
correspond to the visual clusters in Figure~\ref{fig:overview} -- Louvain recovers the political
divide without ever seeing the labels. The remaining communities are small (2--35 blogs) and capture
peripheral or boundary blogs.

\begin{table}[h]
\centering
\caption{Largest detected communities, original graph}
\label{tab:profile}
\small
\begin{tabular}{@{}rrlr@{}}
\toprule
Size & Purity & Majority & Top node (PageRank) \\
\midrule
649 & 0.938 & conservative & instapundit.com \\
518 & 0.977 & liberal & dailykos.com \\
35 & 0.800 & liberal & etalkinghead.com \\
4 & 1.000 & conservative & batr.net \\
3 & 1.000 & liberal & tinylittledots.com \\
\bottomrule
\end{tabular}
\end{table}

\section{Community Evaluation}
\label{sec:eval}

\subsection{Pairwise confusion matrix}
Because detected community IDs are arbitrary and their count need not equal 2, comparison is done
over \emph{unordered node pairs}: TP = same detected community \& same orientation; TN = different
community \& different orientation; FP = same community \& different orientation; FN = different
community \& same orientation. This sidesteps the label-permutation problem completely, since
renumbering communities cannot change which pairs are same/different -- a direct ID-to-label match
is not even well defined here, as there are 11 detected communities but only 2 ground-truth labels.
Table~\ref{tab:conf} reports the declared-run (seed 0) matrix and external metrics for all three
graphs.

\begin{table}[h]
\centering
\caption{Declared-run (seed 0) pairwise confusion matrix and metrics}
\label{tab:conf}
\small
\begin{tabular}{@{}lrrrr@{}}
\toprule
 & Original & Targeted & Random \\
\midrule
TP & 314{,}163 & 194{,}153 & 292{,}054 \\
TN & 342{,}068 & 329{,}230 & 324{,}977 \\
FP & 30{,}628 & 21{,}124 & 25{,}209 \\
FN & 59{,}172 & 157{,}013 & 59{,}280 \\
Accuracy & 0.880 & 0.746 & 0.880 \\
Precision & 0.911 & 0.902 & 0.921 \\
Recall & 0.842 & 0.553 & 0.831 \\
F1 & 0.875 & 0.686 & 0.874 \\
Modularity & 0.426 & 0.433 & 0.424 \\
\# communities & 11 & 61 & 14 \\
\bottomrule
\end{tabular}
\end{table}

\subsection{Zero-denominator handling}
None of the three declared runs produce a zero denominator for any metric: every graph contains at
least one same-community pair and one same-label pair. Had a run produced a single giant community
or a fully-singleton partition, the corresponding metric would be reported as \emph{undefined}
rather than silently replaced with zero (implemented via a division helper that returns \texttt{None}
on a zero denominator).

\subsection{Discussion}
Precision above 0.90 in every condition means that when the algorithm groups two blogs together, it
is very rarely wrong about their political side. Recall is the metric most sensitive to
fragmentation: it measures how completely same-orientation blogs are kept together, and it is
exactly what collapses under targeted broker removal. A structurally strong community (high internal
density, high modularity) can still disagree with political orientation because Louvain optimises
\emph{link density}, not \emph{label agreement} -- a dense but politically mixed cluster of blogs
that all cite each other on a shared sub-topic would score well on modularity while scoring poorly on
precision/recall, and the small mixed community around etalkinghead.com (80\% liberal, Table
\ref{tab:profile}) is a mild example of exactly this.

\section{Reproducibility and Conclusion}
Every source of randomness (the random-node sample and all five Louvain seeds) is controlled by a
single seed (42), set once and reused throughout; the accompanying notebook runs top to bottom
without manual intervention and writes every table and figure as it executes.

The polblogs network shows strong political homophily, visible directly in its hyperlink structure
and recovered by an unsupervised community-detection procedure with 88\% pairwise accuracy against
the true political labels. PageRank and betweenness are strongly but imperfectly correlated: PageRank
finds the most-cited authorities within each side, while betweenness finds a smaller set of
structural bridges, led by blogsforbush.com. The removal experiment confirms this distinction is not
merely definitional -- deleting the top brokers shatters the community structure roughly six-fold,
while deleting the same number of random blogs does not, showing that influence and brokerage are
functionally different roles in this network and that brokers matter disproportionately to how its
community structure holds together.

\subsection*{Limitations}
The data is a single crawl from February 2005 around a polarised U.S. election, which likely
inflates apparent homophily relative to a typical period and does not generalise to other countries,
time periods, or multi-party systems. A hyperlink is treated as a signal of attention, not agreement:
blogs sometimes link to opponents specifically to criticise them, so some cross-cutting edges the
model treats as evidence of bridging may reflect adversarial engagement rather than genuine
brokerage. Finally, Louvain's stochastic node ordering means exact community counts and boundaries
vary somewhat run to run, which is why five seeds are reported rather than a single number.

\begin{thebibliography}{9}
\bibitem{adamic2005} L. A. Adamic and N. Glance, ``The political blogosphere and the 2004 U.S.
election: divided they blog,'' in \emph{Proc. 3rd Int. Workshop on Link Discovery}, 2005.
\bibitem{blondel2008} V. D. Blondel, J.-L. Guillaume, R. Lambiotte, and E. Lefebvre, ``Fast unfolding
of communities in large networks,'' \emph{J. Stat. Mech.}, 2008.
\bibitem{brin1998} S. Brin and L. Page, ``The anatomy of a large-scale hypertextual web search
engine,'' \emph{Computer Networks and ISDN Systems}, vol. 30, 1998.
\bibitem{freeman1977} L. C. Freeman, ``A set of measures of centrality based on betweenness,''
\emph{Sociometry}, vol. 40, no. 1, 1977.
\end{thebibliography}

\end{document}
